\documentclass[11pt]{article}
\usepackage[utf8]{inputenc}
\usepackage[T1]{fontenc}
\usepackage{amsmath}
\usepackage{amssymb}
\usepackage{multicol}
\usepackage{geometry}
\usepackage{tikz}
\usepackage{enumitem}
\usepackage{xcolor}
\usepackage{titlesec}

% Configurações de layout
\geometry{a4paper, left=1cm, right=1cm, top=0.5cm, bottom=1.2cm}
\setlength{\columnseprule}{0.4pt}

% Cores para títulos
\titleformat{\section}{\normalfont\Large\bfseries\color{blue}}{\thesection}{1em}{}
\titleformat{\subsection}{\normalfont\large\bfseries\color{red}}{\thesubsection}{1em}{}


% Comandos personalizados
\newcommand{\respostaNum}[1]{\textcolor{blue}{\boxed{#1}}}
\newcommand{\respostaTxt}[1]{\textcolor{blue}{#1}}
\newcommand{\resolucao}{\textbf{\textcolor{blue}{Resolução:}}}
\newcommand{\etapa}[1]{\textbf{\textcolor{red}{Etapa #1:}}}

\title{\textcolor{red}{Gabarito Comentado: Média, Mediana e Moda}}
\author{Professor: Jefferson}
\date{}

\begin{document}

\maketitle
\vspace{-0.9cm}

\begin{multicols}{2}

\section*{Instruções}
\textbf{Copie e responda no caderno. As resoluções das questões estão em azul.}


\begin{enumerate}[leftmargin=*]

\item \textbf{Calcule a média dos conjuntos:}

\begin{enumerate}[label=\alph*)]
\item 10, 12, 14, 16, 18

\resolucao
\[
\bar{x} = \frac{10+12+14+16+18}{5} = \frac{70}{5} = \respostaNum{14}
\]

\item 5, 8, 11, 14, 17, 20

\resolucao
\[
\bar{x} = \frac{5+8+11+14+17+20}{6} = \frac{75}{6} = \respostaNum{12,5}
\]

\item 2, 4, 6, 8, 10, 12, 14

\resolucao
\[
\bar{x} = \frac{2+4+6+8+10+12+14}{7} = \frac{56}{7} = \respostaNum{8}
\]

\item 15, 20, 25, 30, 35

\resolucao
\[
\bar{x} = \frac{15+20+25+30+35}{5} = \frac{125}{5} = \respostaNum{25}
\]
\end{enumerate}

\item \textbf{Determine a mediana:}

\begin{enumerate}[label=\alph*)]
\item 3, 5, 7, 9, 11

\resolucao
\textbf{Ordem:} já está ordenado. \\
\textbf{n = 5} (ímpar) → termo central: 3º elemento \\
\respostaNum{7}

\item 2, 4, 6, 8, 10, 12

\resolucao
\textbf{Ordem:} já está ordenado. \\
\textbf{n = 6} (par) → média dos 3º e 4º elementos \\
$Md = \frac{6+8}{2} = \frac{14}{2} = \respostaNum{7}$

\item 10, 15, 20, 25, 30, 35, 40

\resolucao
\textbf{n = 7} (ímpar) → 4º elemento \\
\respostaNum{25}

\item 1, 3, 5, 7, 9, 11, 13, 15

\resolucao
\textbf{n = 8} (par) → média dos 4º e 5º elementos \\
$Md = \frac{7+9}{2} = \frac{16}{2} = \respostaNum{8}$
\end{enumerate}

\item \textbf{Identifique a moda:}

\begin{enumerate}[label=\alph*)]
\item 2, 3, 4, 4, 5, 6, 6, 6, 7

\resolucao

\textbf{Frequências:} 2(1x), 3(1x), 4(2x), 5(1x), 6(3x), 7(1x) \\
Maior frequência: 6 (aparece 3 vezes) \\
\respostaNum{6} (unimodal)

\item 10, 12, 12, 14, 14, 16, 18

\resolucao

\textbf{Frequências:} 10(1x), 12(2x), 14(2x), 16(1x), 18(1x) \\
Dois valores com frequência 2: 12 e 14 \\
\respostaNum{12 \text{ e } 14} (bimodal)

\item 5, 5, 5, 8, 8, 10, 10, 10, 12

\resolucao

\textbf{Frequências:} 5(3x), 8(2x), 10(3x), 12(1x) \\
Dois valores com frequência 3: 5 e 10 \\
\respostaNum{5 \text{ e } 10} (bimodal)

\item 1, 2, 3, 4, 5, 6

\resolucao
Todos os valores aparecem uma única vez \\
\respostaNum{\text{amodal (não tem moda)}}
\end{enumerate}

\item \textbf{Para o conjunto: 12, 15, 18, 21, 24, 27, 30}

\begin{enumerate}[label=\alph*)]
\item Média

\resolucao
\[
\bar{x} = \frac{12+15+18+21+24+27+30}{7} = \frac{147}{7} = \respostaNum{21}
\]

\item Mediana

\resolucao
\textbf{n = 7} (ímpar) → 4º elemento \\
\respostaNum{21}

\item Moda

\resolucao
Todos os valores são diferentes \\
\respostaNum{\text{amodal}}
\end{enumerate}

\item \textbf{Um aluno tirou as seguintes notas: 7,0; 8,5; 6,0; 9,0; 7,5. Qual sua média?}

\resolucao
\[
\bar{x} = \frac{7,0+8,5+6,0+9,0+7,5}{5} = \frac{38,0}{5} = \respostaNum{7,6}
\]

\item \textbf{As idades de um grupo são: 12, 14, 13, 15, 12, 16, 14, 13, 12. Determine:}

\resolucao

\textbf{Ordenando:} 12, 12, 12, 13, 13, 14, 14, 15, 16

\begin{enumerate}[label=\alph*)]
\item Média
\[
    \bar{x} = \frac{12+12+12+13+13+14+14+15+16}{9} \] 
\[= \frac{121}{9} \approx \respostaNum{13,44} \]

\item Mediana
\textbf{n = 9} (ímpar) → 5º elemento = \respostaNum{13}

\item Moda
12 aparece 3 vezes (maior frequência) \\
\respostaNum{12}
\end{enumerate}

\end{enumerate}

\begin{enumerate}[leftmargin=*, start=7]

\item \textbf{Salários: R\$ 1800, 2000, 2200, 2500, 5000}

\begin{enumerate}[label=\alph*)]
\item Média
\[
    \bar{x} = \frac{1800+2000+2200+2500+5000}{5} \]
\[= \frac{13500}{5} = \respostaNum{R\$ 2700} \]

\item Mediana

\textbf{Ordenando:} 1800, 2000, 2200, 2500, 5000 \\
n = 5 → 3º elemento = \respostaNum{R\$ 2200}

\item Melhor representação

\resolucao
A \respostaTxt{mediana (R\$ 2200)} representa melhor, pois a média (R\$ 2700) é puxada para cima pelo salário extremo de R\$ 5000.
\end{enumerate}

\item \textbf{Notas: 5, 6, 7, 8, 8, 9, 10, 10, 10}

\begin{enumerate}[label=\alph*)]
\item Média
\[\bar{x} = \frac{5+6+7+8+8+9+10+10+10}{9} = \]
\[\frac{73}{9} \approx \respostaNum{8,11} \]

\item Mediana
n = 9 → 5º elemento = \respostaNum{8}

\item Moda
10 aparece 3 vezes → \respostaNum{10}
\end{enumerate}

\item \textbf{Complete a tabela:}

\begin{center}
    \small
\begin{tabular}{|c|c|c|c|}
\hline
\textbf{Dados} & \textbf{Média} & \textbf{Mediana} & \textbf{Moda} \\
\hline
4, 8, 12, 16, 20 & \textcolor{blue}{12} & \textcolor{blue}{12} & \textcolor{blue}{amodal} \\
\hline
3, 5, 5, 7, 9, 11 & \textcolor{blue}{6,67} & \textcolor{blue}{6} & \textcolor{blue}{5} \\
\hline
10, 20, 30, 40, 50, 60 & \textcolor{blue}{35} & \textcolor{blue}{35} & \textcolor{blue}{amodal} \\
\hline
\end{tabular}
\end{center}

\resolucao

\textbf{1ª linha:} Média = $(4+8+12+16+20)/5 = 60/5 = 12$; Mediana = 12 (3º elemento); todos diferentes → amodal

\textbf{2ª linha:} Média = $(3+5+5+7+9+11)/6 = 40/6 \approx 6,67$; Mediana = $(5+7)/2 = 6$; Moda = 5

\textbf{3ª linha:} Média = $(10+20+30+40+50+60)/6 = 210/6 = 35$; Mediana = $(30+40)/2 = 35$; amodal

\item \textbf{Temperaturas: 22, 24, 23, 25, 24, 26, 24}

\resolucao

\textbf{Ordenando:} 22, 23, 24, 24, 24, 25, 26

\begin{enumerate}[label=\alph*)]
\item Média
\[\bar{x} = \frac{22+23+24+24+24+25+26}{7} \] 
\[= \frac{168}{7} = \respostaNum{24^{\circ}C} \]

\item Mediana
n = 7 → 4º elemento = \respostaNum{24^{\circ}C}

\item Moda
24 aparece 3 vezes → \respostaNum{24^{\circ}C}
\end{enumerate}

\item \textbf{Pesquisa sobre número de irmãos:}

\[
\begin{array}{c|c}
\text{Nº de irmãos} & \text{Frequência} \\
\hline
0 & 5 \\
1 & 8 \\
2 & 6 \\
3 & 3 \\
4 & 2 \\
\end{array}
\]

\resolucao

\textbf{Total de pessoas:} $5+8+6+3+2 = 24$

\begin{enumerate}[label=\alph*)]
\item Média
\[\bar{x} = \frac{0\cdot5 + 1\cdot8 + 2\cdot6 + 3\cdot3 + 4\cdot2}{24} \]
\[= \frac{0+8+12+9+8}{24} = \frac{37}{24} \approx \respostaNum{1,54} \]

\item Mediana
24 elementos (par) → média dos 12º e 13º elementos \\
Ordenando: cinco 0's (posições 1-5), oito 1's (posições 6-13) \\
12º e 13º são ambos 1 → $Md = \frac{1+1}{2} = \respostaNum{1}$

\item Moda
Maior frequência: 1 (8 pessoas) → \respostaNum{1}
\end{enumerate}

\item \textbf{Vendas: 12, 15, 10, 18, 20, 15, 14}

\resolucao

\textbf{Ordenando:} 10, 12, 14, 15, 15, 18, 20

\begin{enumerate}[label=\alph*)]
\item Média
\[\bar{x} = \frac{10+12+14+15+15+18+20}{7} \]
\[= \frac{104}{7} \approx \respostaNum{14,86} \]

\item Mediana
n = 7 → 4º elemento = \respostaNum{15}

\item Moda
15 aparece 2 vezes → \respostaNum{15}
\end{enumerate}

\end{enumerate}

\end{multicols}

\end{document}
